\chapter{Moduł I: Silnik fizyczny symulacji lotów rakietowych.}
\label{ch:silnik-fizyczny-symulacji-lotow-rakietowych.}

Silnik fizyczny do symulacji lotów rakietowych to pierwszy z modułów
niniejszej pracy inżynierskiej.
Za pomocą tego projektu osiągnięty zostaje
cel pośredni pracy - stworzenie środowiska do uczenia algorytmów.
Autorem
tego modułu jest Jakub Kindlik.\\
W tym rozdziale przedstawione zostały teoretyczne podstawy fizyczne, użyte
technologie, oraz
trzy główne komponenty
tego modułu:

\begin{itemize}
    \item Projekt silnika fizycznego.
    \item Model sił działających na rakietę.
    \item Moduł wizualizacyjny symulacji.
\end{itemize}

Omówione zostały także zagadnienia związane ze
sprawdzeniem dokładności symulacji oraz ewolucją projektu.\\

\section{Założenia projektowe.}\label{sec:zalozenia-projektowe.}
W trakcie pracy nad modułem silnika fizycznego symulacji lotów rakietowych poczyniono szereg założeń, które miały na celu ułatwienie implementacji oraz zapewnienie odpowiedniej wydajności symulacji. Poniżej przedstawiono najważniejsze z nich:

\begin{itemize}

\item \textbf{Modelowanie uproszczonych sił działających na rakietę.} \\
W celu zmniejszenia złożoności obliczeniowej model sił działających na rakietę zostanie odpowiednio uproszczony. Uwzględnione zostaną jedynie kluczowe siły, takie jak grawitacja, opór powietrza i siła ciągu silnika. Takie podejście umożliwia łatwiejsze zamodelowanie systemu i redukuje zapotrzebowanie na zasoby obliczeniowe.

\item \textbf{Optymalizacja pod kątem pracy na komputerze osobistym.} \\
Symulacje będą wykonywane na komputerze osobistym, dlatego silnik fizyczny
musi być zoptymalizowany pod kątem wydajności, jak i okrojony w kwestii
skomplikowanych obliczeń. Ograniczenie
zasobożerności
jest konieczne, aby zapewnić płynne działanie nawet na sprzęcie o przeciętnych parametrach.

\item \textbf{Uproszczona geometria rakiety.} \\
W celu ograniczenia złożoności modelu fizycznego rakieta będzie miała kształt walca i zostanie pozbawiona skrzydeł. Taka decyzja pozwala na skupienie się na kluczowych aspektach dynamiki lotu bez nadmiernego komplikowania obliczeń.

\item \textbf{Sterowanie poprzez obracany silnik.} \\
Układ sterujący rakietą będzie się opierał na możliwości obracania silnika,
co bezpośrednio wpłynie na trajektorię lotu. Taki model jest prosty do
zamodelowania, a jednocześnie pozwala na sterowanie rakietą przez algorytmy
samouczące, poprzez zmianę jednego parametru.

\item \textbf{Możliwość modyfikacji parametrów symulacji w czasie rzeczywistym.} \\
Symulacja zostanie zaprojektowana w sposób umożliwiający ingerencję w jej
parametry w trakcie trwania. W szczególności modyfikowanym parametrem
powinien być kąt nachylenia silnika. Taka funkcjonalność pozwala na
uczenie algorytmów sterujących w czasie rzeczywistym.

\item \textbf{Możliwość wizualizacji symulacji.} \\
Symulacja będzie wyposażona w opcjonalną funkcję wizualizacji, co umożliwi analizę trajektorii i zachowania rakiety w sposób intuicyjny. Jednak wizualizacja nie będzie wymagana, co pozwoli ograniczyć zużycie zasobów obliczeniowych w trakcie symulacji.

\end{itemize}

Założenia te konieczne są do stworzenia silnika fizycznego, który będzie w
stanie współpracować z modułem algorytmu sterującego rakietą, jednocześnie
będąc na tyle prostym, aby nie obciążać zasobów komputera osobistego.


\section{Teoretyczne podstawy fizyczne.}\label{sec:teoretyczne-podstawy.}

Tworzenie silnika symulacji lotów rakietowych wymaga zrozumienia podstawowych
zasad fizyki odnoszących się do ruchu rakiety w atmosferze.


\section{Użyte technologie.}\label{sec:uzyte-technologie.}

Do stworzenia silnika fizycznego symulacji lotów rakietowych wykorzystane zostały następujące technologie:

\begin{itemize}
    \item \textbf{Python 3.13} - język programowania
    \item zestaw bibliotek
    \begin{itemize}
        \item bibioteka1....
    \end{itemize}
\end{itemize}

\section{Projekt silnika fizycznego.}\label{sec:projekt-silnika-fizycznego.}
\section{Model sił działających na rakietę.}\label{sec:model-si-dziaajacych-na-rakiete.}
\section{Moduł wizualizacyjny symulacji.}\label{sec:modul-wizualizacyjny-symulacji.}
\section{Dokładność symulacji.}\label{sec:dokadnosc-symulacji}
\section{Ewolucja projektu i wnioski.}\label{sec:ewolucja-projektu-i-wnioski}